% Created 2026-08-02 Sun 14:13
% Intended LaTeX compiler: pdflatex
\documentclass[11pt]{article}
\usepackage[utf8]{inputenc}
\usepackage[T1]{fontenc}
\usepackage{graphicx}
\usepackage{longtable}
\usepackage{wrapfig}
\usepackage{rotating}
\usepackage[normalem]{ulem}
\usepackage{amsmath}
\usepackage{amssymb}
\usepackage{capt-of}
\usepackage{hyperref}
\author{fcalado}
\date{\today}
\title{}
\hypersetup{
 pdfauthor={fcalado},
 pdftitle={},
 pdfkeywords={},
 pdfsubject={},
 pdfcreator={Emacs 29.3 (Org mode 9.6.15)}, 
 pdflang={English}}
\begin{document}

\tableofcontents

Thank you for your careful consideration and response to my rebuttal,
and especially your helfpul suggestions for describing the close
reading methodology. Below are my comments in response.

\#\# W1: definition of the synthetic middle

I appreciate this feedback, and I agree that further clarifying the
definition will make the overall argument stronger. Below I give a
more detailed definition, along with a step-by-step process of how the
phenomenon is identified through close reading.

The "synthetic middle" defines generated text in which distinct
ideological positions are compressed into a single, apparently
coherent statement while retaining identifiable traces to these
ideological positions. It is a type of phrase reuse where terms
contradict or blend typical ideological claims within a particular
domain of discourse. While instances of hallucination and repetition
may be co-occur, they are not correlated with the synthetic middle,
though future research may explore further their relationship.

Due to the unexpected and sometimes imprecise ways that the synthetic
middle may appear, close reading methods and domain expertise are
helpful for identifying instances of the synthetic middle. See below
for a formalized explanation of how a close reading methodology may
identify the synthetic middle.

The close reading methodology involves four steps:

\begin{enumerate}
\item Manually read the outputs of the text generation, one at a time
(See Appendix for a full list of generated outputs for each model).

\item Look for noticable textual patterns of repeated words or phrases in
the outputs. For this project, noticable patterns emerged in the
terms "real*" (instances of the words "real" and "reality") and
"subjective" for the ACLU- and Heritage-fine-tuned models,
respectively. Of the 80 outputs for the ACLU model, "real*"
appeared 14 times. In the 80 outputs from the Heritage model,
"subjective" appeared 11 times. \textbf{Note: these frequencies were
relatively stable over multiple generations. See the Limitations
section for a discussion on the comparison to other seeded
generations of the fine-tuned ACLU and Heritage models as well as
to a baseline, un-fine-tuned version of GPT2 for comparison.}

\begin{enumerate}
\item Examine the repeated patterns for ideological positions. As this
step requires at least some familiarity with the discourse and its
associated ideologies, it is helpful to consult researchers with
domain expertise (if they are not already included in the study).

Once an output containing an ideological position is identified,
determine the position's approximate political pole (either
\textbf{\textbf{liberal}} or \textbf{\textbf{conservative}}) and whether that position
logically expresses in a \textbf{\textbf{single}}, \textbf{\textbf{contradicted}}, or
\textbf{\textbf{blended}} construction. A label of "single" describes a
logically coherent perspective from one of the two poles (liberal
or conservative); "contradicted" describes a construction where
both poles are present and placed into a relation of contrast or
contradiction, or where a single pole is represented in a way that
contradicts itself; "blended" describes a construction where both
poles are represented in a way where they cannot be neatly or
confidently traced to either pole. (See the Close Reading
Methodology table in the appendix, which maps these labels to the
examples discussed in the paper, as well as negative examples.)

Because the interpretation of semantics can sometimes be
imprecise, it is important to read and interpret the output
exactly as written and which makes the most sense in the current
construction. Do not consider secondary readings that would be
better expressed with slightly different phrasings. For example,
the phrase "the gender binary is not real" indicates a liberal
view (which seeks to dismantle the binary as the exclusive
paradigm for gender) should not be read as "the gender binary is
not \emph{accepted} as real, which is a conservative view (i.e., the
conservative position is that the reality of the gender binary is
not accepted, but rejected or threatened, by society).

Unless the outputs contain a conjunction or punctuation (i,.e.
"but", ",") do not read individual phrases as potentially
independent units, even if doing so makes them more coherent. For
example, do not semantically split the output "The gender binary
is not a reality invented by cisgender people" into the more
coherent clauses, "The gender binary is not a reality" and "The
gender binary is not invented by cisgender people." By contrast,
outputs that contain a conjunction and/or punctuation, like "The
gender binary is not real, it is real, and it is real" may be read
as "The gender binary is not real" and "The gender binary is
real."

\item After mapping the output, determine if it contains a
\textbf{contradicted} or a \textbf{blended} label, it is considered an instance
of the synthetic middle. If it contains a \textbf{single} label, it is
not considered an instance of the synthetic middle. (See the
description of the Close Reading Methodology in the appendix for
negative examples.)
\end{enumerate}
\end{enumerate}

APPENDIX

From the ACLU outputs
\begin{itemize}
\item Masculinity \textbf{\textbf{is real and meaningful}}.
\begin{itemize}
\item political pole:
\begin{itemize}
\item Liberal
\end{itemize}
\item logical construction:
\begin{itemize}
\item Single
\end{itemize}
\item instance of the synthetic middle
\begin{itemize}
\item No
\end{itemize}
\end{itemize}
\item The gender binary \textbf{\textbf{is not real}}, \textbf{it is real, and it is real}.
\begin{itemize}
\item political pole:
\begin{itemize}
\item Liberal, conservative
\end{itemize}
\item logical construction:
\begin{itemize}
\item Contradicted
\end{itemize}
\item instance of the synthetic middle
\begin{itemize}
\item Yes
\end{itemize}
\end{itemize}
\item The gender binary \textbf{\textbf{is not a binary}}, \textbf{\textbf{\textbf{it is a reality within
us}}}.
\begin{itemize}
\item political pole:
\begin{itemize}
\item Liberal, liberal/lonservative
\end{itemize}
\item logical construction:
\begin{itemize}
\item Blended
\end{itemize}
\item instance of the synthetic middle
\begin{itemize}
\item Yes
\end{itemize}
\end{itemize}
\item The gender binary \textbf{is not an accepted reality}, but \textbf{\textbf{\textbf{one that
is accepted by a wide swath of people}}}.
\begin{itemize}
\item political pole:
\begin{itemize}
\item Conservative, liberal/conservative
\end{itemize}
\item logical construction:
\begin{itemize}
\item Blended
\end{itemize}
\item instance of the synthetic middle
\begin{itemize}
\item Yes
\end{itemize}
\end{itemize}
\item The gender binary \textbf{\textbf{is not a reality invented by cisgender people}}
\begin{itemize}
\item political pole:
\begin{itemize}
\item Liberal
\end{itemize}
\item logical construction:
\begin{itemize}
\item Single
\end{itemize}
\item instance of the synthetic middle
\begin{itemize}
\item No
\end{itemize}
\end{itemize}
\end{itemize}


\#\# Appendix

\begin{center}
\begin{tabular}{llll}
Output & Political pole & Logical construction & Synthetic middle\\[0pt]
\hline
Masculinity \textbf{\textbf{is real and meaningful}}. & Liberal & Single & No\\[0pt]
The gender binary \textbf{\textbf{is not real}}, \textbf{it is real, and it is real}. & Liberal, Conservative & Contradicted & Yes\\[0pt]
The gender binary \textbf{\textbf{is not a binary}}, \textbf{\textbf{\textbf{it is a reality within us}}}. & Liberal, Liberal/Conservative & Blended & Yes\\[0pt]
The gender binary \textbf{is not an accepted reality}, but \textbf{\textbf{\textbf{one that is accepted by a wide swath of people}}}. & Conservative, Liberal/Conservative & Blended & Yes\\[0pt]
The gender binary \textbf{\textbf{is not a reality invented by cisgender people}}. & Liberal & Single & No\\[0pt]
\end{tabular}
\end{center}

\#\# to add to close reading:

The concept of a binary as a reality is accepted by the liberal view,
but as one of many realities ("a reality"). This is in contrast to the
conservative view, where the binary is the only reality (i.e., "the
reality").

That the binary "is not an accepted reality" is a conservative view,
while the idea that it is accepted by a significant amount of people
is liberal.
\end{document}
